
\subsubsection{Quality-Based Data Curation}

V-Information provides a principled quality metric using pointwise information content, demonstrating that quality-based reduction maintains performance at small scales. DATAMASK made qualitative observations that quality metrics show diminishing returns at large scale but did not quantify trajectories or test sequential ordering effects. We extend this work by measuring quality saturation as a continuous function (regression slope $-2.62pp/log-scale$, p=0.008) and testing compositional interactions.

\subsubsection{Diversity-Based Data Curation}

Feature Activation Coverage (FAC) demonstrated $\rho$=0.90 correlation with downstream performance on large-scale datasets. Prior work validated FAC effectiveness across domains but did not analyze scale-dependent trends or interaction with quality filtering. We model diversity persistence trajectories from literature and test whether diversity effectiveness depends on operating on quality-filtered versus raw data.

\subsubsection{Data Mixing and Composition}

Data Mixing Laws established that small proxy models predict large-scale mixture performance when combining domain proportions, validating our use of GPT-2 Small for experiments. However, Mixing Laws assume sources mix independently with additive contributions---an order-invariance assumption that holds for domain mixing. We extend this framework by testing whether order matters for sequential filtering steps within a single source, finding large effect sizes (d=0.76--2.48) that indicate order-invariance does not hold for quality-diversity filtering sequences.

RegMix optimizes mixture proportions via regression but similarly assumes order-invariant effects. Our findings suggest that optimization frameworks could potentially improve by incorporating sequential ordering: apply quality filters before diversity sampling within each source, then optimize domain proportions.

\subsubsection{Positioning}

We extend prior work by (1) quantifying saturation and persistence as scale-dependent trajectories rather than binary observations, (2) demonstrating compositional ordering effects in sequential filtering steps, and (3) providing prescriptive guidance grounded in controlled experiments showing quality-first universality.
